\documentclass[aspectratio=169,10pt]{beamer}
\usetheme[numbering=fraction,progressbar=frametitle]{metropolis}
\usepackage{fontspec}
\setsansfont{Latin Modern Sans}
\usepackage{amsmath,amssymb,mathtools,graphicx,booktabs}
\usepackage{polyglossia}
\setdefaultlanguage{english}
\setbeamertemplate{navigation symbols}{}
\setbeamerfont{frametitle}{size=\large}
\setbeamerfont{normal text}{size=\normalsize}
\setlength{\parskip}{.45em}
\hypersetup{pdfauthor={},pdfsubject={Lecture slides},pdfkeywords={},pdftitle={Functions: representations and reference functions}}
\providecommand{\N}{\mathbb{N}}
\providecommand{\Z}{\mathbb{Z}}
\providecommand{\D}{\mathbb{D}}
\providecommand{\Q}{\mathbb{Q}}
\providecommand{\R}{\mathbb{R}}
\providecommand{\C}{\mathbb{C}}
\providecommand{\vect}[1]{\overrightarrow{#1}}
\providecommand{\vi}{\vec{\imath}}
\providecommand{\vj}{\vec{\jmath}}
\providecommand{\vk}{\vec{k}}
\providecommand{\norm}[1]{\left\lVert #1\right\rVert}
\providecommand{\vc}[2]{\begin{pmatrix}#1\\#2\end{pmatrix}}
\providecommand{\deter}[2]{\det\left(#1,#2\right)}
\providecommand{\abs}[1]{\left\lvert #1\right\rvert}
\providecommand{\dd}{\,\mathrm{d}}
\providecommand{\e}{\mathrm{e}}
\providecommand{\Card}{\operatorname{Card}}
\providecommand{\E}{\operatorname{E}}
\providecommand{\V}{\operatorname{V}}
\providecommand{\degre}{^{\circ}}
\providecommand{\un}[1]{\,\mathrm{#1}}
\providecommand{\sci}[2]{#1\times10^{#2}}
\providecommand{\pgcd}{\operatorname{PGCD}}
\providecommand{\ppcm}{\operatorname{PPCM}}
\title{Functions: representations and reference functions}
\subtitle{Lecture and worked examples}
\date{}
\author{}
\begin{document}
\begin{frame}[plain]
\vfill
{\usebeamerfont{title}\inserttitle\par}
\medskip
{\usebeamerfont{subtitle}\insertsubtitle\par}
\medskip
{\usebeamercolor[fg]{progress bar}\rule{\textwidth}{.5pt}}
\vfill
\end{frame}
\begin{frame}{Function, image and preimage}
A real-valued function associates to each real number \(x\) of its domain \(D_f\) a unique real number \(f(x)\), called the image of \(x\). A preimage of \(y\) is a real number \(x\in D_f\) such that \(f(x)=y\). We write \(f:x\mapsto f(x)\). A number may have no preimage, one, or several.
\end{frame}
\begin{frame}{Function, image and preimage}
Let \(g(x)=x^2-4x+32\) on \(\R\). Calculate \(g(-5)\). Let \(h(x)=3x+5\); determine the preimage of \(6\) under \(h\).
\end{frame}
\begin{frame}{Worked answer: Function, image and preimage}
\(g(-5)=25+20+32=77\): \(77\) is the image of \(-5\), and \(-5\) is a preimage of \(77\). For \(h\), \(3x+5=6\iff x=1/3\). The only preimage is \(1/3\).
\end{frame}
\begin{frame}{Function, image and preimage}
For the square function, calculate the images of \(-3\) and \(3\), then the preimages of \(9\), \(0\) and \(-1\).
\end{frame}
\begin{frame}{Worked answer: Function, image and preimage}
\(f(-3)=f(3)=9\). The preimages of \(9\) are \(-3,3\); \(0\) has a single preimage, \(0\); \(-1\) has none in \(\R\).
\end{frame}
\begin{frame}[fragile]{Python}
\small
\begin{verbatim}
def g(x):
    return x*x - 4*x + 32
\end{verbatim}
\end{frame}
\begin{frame}{Function, image and preimage}
Distinguish \texttt{g(-5)} from \texttt{x = -5}.
\end{frame}
\begin{frame}{Worked answer: Function, image and preimage}
The first calls the function and returns \(77\). The second assigns \(-5\) to the variable \texttt{x}, without calculating the image under \(g\).
\end{frame}
\begin{frame}{Domain of definition}
\(D_f\) contains the values for which the process defining \(f\) makes sense. A denominator must be nonzero; a number under a square root must be positive or zero. Also respect any domain imposed by the statement or the context.
\end{frame}
\begin{frame}{Domain of definition}
Determine the domains of definition:
\[f(x)=3x^2-5x+2,\quad g(x)=\dfrac{2x+1}{x-4},\quad h(x)=\dfrac{x}{x^2-9},\quad k(x)=\sqrt{2x-6}.\]
\end{frame}
\begin{frame}{Worked answer: Domain of definition}
\(D_f=\R\). \(D_g=(-\infty,4)\cup(4,+\infty)\) since \(x-4\ne0\). \(D_h=(-\infty,-3)\cup(-3,3)\cup(3,+\infty)\) since \(x^2-9\ne0\). \(D_k=[3,+\infty)\) since \(2x-6\geq0\).
\end{frame}
\begin{frame}{Domain of definition}
The square root of a negative number does not exist in \(\R\); it is not “a negative square root” that is excluded.
\end{frame}
\begin{frame}{Domain of definition}
A rectangle of perimeter \(20\) m has width \(x\) m. Express its area and state the domain if both dimensions must be strictly positive.
\end{frame}
\begin{frame}{Worked answer: Domain of definition}
The length is \(10-x\), so \(A(x)=x(10-x)\) in m\(^2\). \(x>0\) and \(10-x>0\) give \(D_A=(0,10)\). If degenerate rectangles were allowed, we could take \([0,10]\).
\end{frame}
\begin{frame}{Table of values and representative curve}
\(\mathcal C_f\) is the set of points \((x,f(x))\), with \(x\in D_f\). Thus \(M(a,b)\in\mathcal C_f\iff a\in D_f\) and \(b=f(a)\).
\end{frame}
\begin{frame}{Table of values and representative curve}
\begin{columns}[T,onlytextwidth]
\begin{column}{.57\textwidth}
On \([-1,4]\), \(f(x)=x^2-2x-1\). Complete the table, then draw the curve over the whole interval.
\end{column}
\begin{column}{.4\textwidth}
\centering
\includegraphics[width=\linewidth,height=4.5cm,keepaspectratio]{figures/seconde-06-1-en.pdf}
\end{column}
\end{columns}
\end{frame}
\begin{frame}{Worked answer: Table of values and representative curve}
Images: \(2\); \(0.25\); \(-1\); \(-2\); \(-1.75\); \(-1\); \(2\); \(7\). Plot the points then draw the parabola. Segments joining only these points would give just an approximation.
\end{frame}
\begin{frame}{Graphical solving of equations and inequalities}
The solutions of \(f(x)=k\) are the abscissae of the intersections of \(\mathcal C_f\) with \(y=k\). For \(f(x)\leq k\), read the abscissae where the curve is below this line, intersections included; with \(<\), the intersections are excluded.
\end{frame}
\begin{frame}{Graphical solving of equations and inequalities}
\begin{columns}[T,onlytextwidth]
\begin{column}{.57\textwidth}
The curve represents \(f(x)=1+\dfrac14(x+3)(x+1)(x-2)\) on \([-4,3]\). Solve graphically \(f(x)=1\), \(f(x)\leq1\) and \(f(x)>1\).
\end{column}
\begin{column}{.4\textwidth}
\centering
\includegraphics[width=\linewidth,height=4.5cm,keepaspectratio]{figures/seconde-06-2-en.pdf}
\end{column}
\end{columns}
\end{frame}
\begin{frame}{Worked answer: Graphical solving of equations and inequalities}
\(f(x)=1\): \(-3,-1,2\). \(f(x)\leq1\): \([-4,-3]\cup[-1,2]\). \(f(x)>1\): \((-3,-1)\cup(2,3]\). The endpoints \(-4,3\) come from the domain. The factorisation checks the intersections exactly; the algebraic study of the sign will be consolidated in sequence 8.
\end{frame}
\begin{frame}{Comparing two functions; reading a sign}
\(f(x)=g(x)\) is read at the intersections of the curves. For \(f(x)<g(x)\), \(\mathcal C_f\) is strictly below \(\mathcal C_g\) at the same abscissa. The sign of \(f(x)\) is read relative to the x-axis.
\end{frame}
\begin{frame}{Comparing two functions; reading a sign}
\begin{columns}[T,onlytextwidth]
\begin{column}{.57\textwidth}
On \([-3,3]\), \(f(x)=x^2\) and \(g(x)=x+2\). Solve \(f=g\), \(f<g\), \(f\geq g\) and determine the sign of \(g\).
\end{column}
\begin{column}{.4\textwidth}
\centering
\includegraphics[width=\linewidth,height=4.5cm,keepaspectratio]{figures/seconde-06-3-en.pdf}
\end{column}
\end{columns}
\end{frame}
\begin{frame}{Worked answer: Comparing two functions; reading a sign}
\(f=g\) for \(x=-1,2\); \(f<g\) on \((-1,2)\); \(f\geq g\) on \([-3,-1]\cup[2,3]\). \(g\) is strictly negative on \([-3,-2)\), zero at \(-2\), strictly positive on \((-2,3]\). Sign and variations are different pieces of information.
\end{frame}
\begin{frame}{Affine functions and absolute value}
An affine function is \(f(x)=mx+p\) on \(\R\). It is linear if \(p=0\), constant if \(m=0\). For \(m\ne0\), it vanishes at \(-p/m\); its sign is that of \(m\) to the right of this number.
The function \(x\mapsto|x|\) is defined on \(\R\), positive or zero, and vanishes only at \(0\).
\end{frame}
\begin{frame}{Affine functions and absolute value}
\begin{columns}[T,onlytextwidth]
\begin{column}{.57\textwidth}
Classify (several answers possible) as affine, linear, constant: \(5x+2\); \(5x\); \(7+2x-7\); \(6\); \(6(4x-2)\); \(6x+5-6x\).
\end{column}
\begin{column}{.4\textwidth}
\centering
\includegraphics[width=\linewidth,height=4.5cm,keepaspectratio]{figures/seconde-06-4-en.pdf}
\end{column}
\end{columns}
\end{frame}
\begin{frame}{Worked answer: Affine functions and absolute value}
All are affine. \(5x\) and \(2x\) are linear; \(6\) and \(5\) are constant; \(24x-12\) is neither linear nor constant.
\end{frame}
\begin{frame}{Affine functions and absolute value}
If \(k<0\), \(|x|=k\) and \(|x|\leq k\) have no solution, and \(|x|\geq k\) is true on \(\R\). If \(k=0\), the first two give \(x=0\), the third is true on \(\R\). For \(k>0\):
\[|x|=k\iff x\in\{-k,k\},\quad |x|\leq k\iff x\in[-k,k],\]
\[|x|\geq k\iff x\in(-\infty,-k]\cup[k,+\infty).\]
\end{frame}
\begin{frame}{Affine functions and absolute value}
Solve \(|x|=5\), \(|x|=-3\), \(|x|\leq3\), \(|x|\geq4\).
\end{frame}
\begin{frame}{Worked answer: Affine functions and absolute value}
\(\{-5,5\}\); \(\varnothing\); \([-3,3]\); \((-\infty,-4]\cup[4,+\infty)\).
\end{frame}
\begin{frame}{Square and inverse functions}
\begin{columns}[T,onlytextwidth]
\begin{column}{.57\textwidth}
\(x\mapsto x^2\) is defined on \(\R\); its curve is a parabola. \(x^2\geq0\), with equality only at \(0\).
\(x\mapsto1/x\) is defined on \(\R^*\); its curve is a hyperbola. \(1/x\) has the sign of \(x\) and never vanishes.
\end{column}
\begin{column}{.4\textwidth}
\centering
\includegraphics[width=\linewidth,height=2.25cm,keepaspectratio]{figures/seconde-06-5-en.pdf}
\includegraphics[width=\linewidth,height=2.25cm,keepaspectratio]{figures/seconde-06-6-en.pdf}
\end{column}
\end{columns}
\end{frame}
\begin{frame}{Square and inverse functions}
If \(k<0\), \(x^2=k\) and \(x^2\leq k\) have no solution; \(x^2\geq k\) is true on \(\R\). If \(k=0\), the first two give \(x=0\), the third is true on \(\R\). For \(k>0\):
\[x^2=k\iff x\in\{-\sqrt k,\sqrt k\},\quad x^2\leq k\iff x\in[-\sqrt k,\sqrt k],\]
\[x^2\geq k\iff x\in(-\infty,-\sqrt k]\cup[\sqrt k,+\infty).\]
\end{frame}
\begin{frame}{Square and inverse functions}
Solve \(x^2=4\), \(x^2=-1\), \(x^2\geq4\), \(x^2\geq-1\), \(1/x=4\), \(1/x=0\).
\end{frame}
\begin{frame}{Worked answer: Square and inverse functions}
\(\{-2,2\}\); \(\varnothing\); \((-\infty,-2]\cup[2,+\infty)\); \(\R\); \(\{1/4\}\); \(\varnothing\). More generally, \(1/x=k\) has the unique solution \(1/k\) if \(k\ne0\), none if \(k=0\).
\end{frame}
\begin{frame}{Inequalities with the inverse function}
Respect \(x\ne0\) and read both branches of the hyperbola. Do not multiply an inequality by \(x\) without knowing its sign.
\end{frame}
\begin{frame}{Inequalities with the inverse function}
Solve \(1/x\geq2\), \(1/x\leq-1\), \(1/x<1/2\) and \(1/x>-2\).
\end{frame}
\begin{frame}{Worked answer: Inequalities with the inverse function}
\((0,1/2]\); \([-1,0)\); \((-\infty,0)\cup(2,+\infty)\); \((-\infty,-1/2)\cup(0,+\infty)\). A strict inequality excludes the intersection with \(y=k\); \(0\) is always forbidden.
\end{frame}
\begin{frame}{Inequalities with the inverse function}
A student writes: “\(1/x\leq2\iff1\leq2x\), so \(x\geq1/2\)”. What is missing?
\end{frame}
\begin{frame}{Worked answer: Inequalities with the inverse function}
They implicitly assumed \(x>0\). For \(x<0\), \(1/x<0<2\): these values also work. The correct set is \((-\infty,0)\cup[1/2,+\infty)\). The number \(-1\) refutes their equivalence.
\end{frame}
\begin{frame}{Using square root and cube}
\begin{columns}[T,onlytextwidth]
\begin{column}{.57\textwidth}
\(x\mapsto\sqrt x\) is defined on \([0,+\infty)\) and takes values that are positive or zero. For \(k\geq0\), \(\sqrt x=k\iff x=k^2\); for \(k<0\), no solution.
\(x\mapsto x^3\) is defined on \(\R\). Every real number \(k\) has a unique preimage, written \(\sqrt[3]{k}\).
\end{column}
\begin{column}{.4\textwidth}
\centering
\includegraphics[width=\linewidth,height=2.25cm,keepaspectratio]{figures/seconde-06-7-en.pdf}
\includegraphics[width=\linewidth,height=2.25cm,keepaspectratio]{figures/seconde-06-8-en.pdf}
\end{column}
\end{columns}
\end{frame}
\begin{frame}{Using square root and cube}
Solve \(\sqrt x=3\); \(\sqrt x=-2\); \(\sqrt x\leq2\); \(\sqrt x>3\); \(x^3=-8\); \(x^3\leq27\).
\end{frame}
\begin{frame}{Worked answer: Using square root and cube}
\(\{9\}\); \(\varnothing\); \([0,4]\); \((9,+\infty)\); \(\{-2\}\); \((-\infty,3]\). Respect the domain for square roots. The cube function preserves order.
\end{frame}
\begin{frame}{Using square root and cube}
Compare \(x^2\) and \(x\) for \(x=1/2\), then \(x=2\). Is the square of a positive number always greater than the number itself?
\end{frame}
\begin{frame}{Worked answer: Using square root and cube}
\(1/4<1/2\) but \(4>2\). The claim is false; \(1/2\) is a counterexample. General study in sequence 10.
\end{frame}
\end{document}
